% Generated by generate.py; do not edit.
\subsection*{Problem}
\begin{description}
\item[Problem] One instance to be scheduled: a set of students, a set of exams, and the number of batches available.
\item[Student] A person who must sit zero or more exams; identified by a unique name.
\item[Exam] A single final exam sat by one or more students; identified by a unique name and carrying a weight.
\item[Weight] The priority of an exam: the cost of leaving it unscheduled. Always positive; 1 when not given, so by default a schedule places as many exams as it can.\par
{\small\itshape Avoid: profit, importance}
\item[Batch] One of a fixed number of exam periods; all exams in the same batch take place simultaneously. The \enquote{slot} in the problem's name, Exam-Slot Assignment, is a batch.\par
{\small\itshape Avoid: slot, time slot, TimeSlot, period}
\end{description}
\subsection*{Conflict Graph}
\begin{description}
\item[Conflict] A relation between two distinct exams that share at least one student, so they must not be in the same batch.\par
{\small\itshape Avoid: contradiction, clash}
\item[Conflict Graph] The exams as nodes and their conflicts as edges, together with the number of batches.\par
{\small\itshape Avoid: course conflict, contradiction matrix}
\end{description}
\subsection*{Solving}
\begin{description}
\item[Solver Configuration] The choices about how a problem is solved rather than what the problem is: the method, its settings, the seed and the penalty coefficient.\par
{\small\itshape Avoid: settings, parameters}
\item[Method] How the compiled formulation is sampled: QAOA or simulated annealing.\par
{\small\itshape Avoid: solver (ambiguous with the Solver Configuration), backend}
\item[Assignment Variable] The binary decision x[Exam, Batch]: 1 when that exam is placed in that batch. pyqubo labels it \texttt{c[i][k]}, as in the notebook.\par
{\small\itshape Avoid: qubit (that is its position in the QAOA circuit)}
\item[Hamiltonian] The energy to minimise, H = \texttt{profit} + A · (\texttt{constraint1} + \texttt{constraint2}), with \texttt{profit} = −Σ weights of scheduled exams; \texttt{constraint1} keeps each exam in at most one batch, \texttt{constraint2} keeps conflicting exams out of the same batch.\par
{\small\itshape Avoid: cost function, objective (for the whole)}
\item[Constraint] One pyqubo \texttt{Constraint} term of the Hamiltonian, labelled in domain terms: \texttt{at-most-one-batch-<exam>}, or \texttt{conflict-<exam>-<exam>-batch-<k>}.\par
{\small\itshape Avoid: penalty term}
\item[Penalty Coefficient] The value fed to the Hamiltonian's placeholder A. Automatic by default, the largest weight plus one; an explicit value must exceed every weight.\par
{\small\itshape Avoid: lambda, penalty weight}
\item[QUBO] The Hamiltonian as pyqubo compiles it: a quadratic polynomial over the assignment variables. It exists only at run time, inside the generated script.\par
{\small\itshape Avoid: QUBO matrix (BQM is pyqubo's name for the same thing)}
\item[QAOA Circuit] The quantum circuit built from the Ising form of the compiled model: initial superposition, p alternating cost and mixer layers, and measurement.\par
{\small\itshape Avoid: ansatz, program, quantum model}
\item[Simulated Annealing] The classical sampler the notebook uses (\texttt{neal}): a number of reads, each a run of a number of sweeps.\par
{\small\itshape Avoid: quantum annealing (the paper's D-Wave hardware, which the chain does not use)}
\end{description}
\subsection*{Solution}
\begin{description}
\item[Schedule] The outcome of solving a problem: the batch of each scheduled exam, as decoded from the lowest-energy sample.\par
{\small\itshape Avoid: plan, solution matrix}
\item[Unscheduled Exam] An exam the schedule places in no batch, accepted at the cost of its weight.\par
{\small\itshape Avoid: missed exam, dropped exam}
\item[Cost {\normalfont(of a schedule)}] The total weight of its unscheduled exams.
\item[Feasible {\normalfont(schedule)}] Breaking no constraint: each exam in at most one batch and no conflict within a batch. Unscheduled exams are allowed. An infeasible schedule lists its broken constraints.\par
{\small\itshape Avoid: valid (the BESSER chain's word for a schedule that places every exam with no conflict: here, feasible with no unscheduled exam)}
\item[Optimal {\normalfont(schedule)}] Feasible, with the lowest cost any feasible schedule of the problem has.
\end{description}
\subsection*{Validation}
\begin{description}
\item[Baseline] Enumerating every schedule (each exam in one batch or none), independently of the QUBO, to find the optimal cost.\par
{\small\itshape Avoid: reference solution, ground truth}
\item[P(optimal)] The fraction of QAOA shots or annealing reads that decode to an optimal schedule.\par
{\small\itshape Avoid: success rate}
\item[P(feasible)] The fraction of QAOA shots or annealing reads that decode to a feasible schedule.\par
{\small\itshape Avoid: success rate}
\end{description}
